% ===================================================================
%  15. TAULA — Merkatua. — Janariak.
%  Traduction 2026 d'apres texto/io, controlee sur texto/fr, texto/en
%  et texto/es. Consigne : voir texto/eu/10-taula-01.tex.
%
%  LE TABLEAU DU MARCHE REND AU BASQUE CE QUE LA GARE ET LE PORT LUI
%  AVAIENT PRIS. Le boudin est « odolkia », le saindoux « urdaia », le
%  fer a repasser « lisaburdina », la lentille « dilista », le hareng
%  « sardinzarra », le pain « ogia ». Ce sont les mots de la maison, et
%  ils sont a la langue.
% ===================================================================

%%K t15-apar-1 apar
\VUsaut{4.0mm}
\VUtitre{15.0pt}{60}{15. TAULA}
\VUsaut{3.0mm}

%%K t15-tit-1 sub
\VUpk{11.4pt}{40}{Merkatua. --- Janariak.}
\VUsaut{3.0mm}

%%K t15-01-1 p
1. --- Behar ditugun janariak hornitzen dizkiguten merkatarien gehiengoa
\VUgras{merkatu estaliaren} \VUgras{teilapean} \textsuperscript{(1)} edo
ondoko kaleetan bildu da.

%%K t15-02-1 p
2. --- \VUgras{Hestebete-saltzailea} \textsuperscript{(2)}, zeinak
\VUgras{urdaiazpikoa} \textsuperscript{(3)}, \VUgras{odolkia}
\textsuperscript{(4)}, \VUgras{lukainka} \textsuperscript{(5)},
\VUgras{saltxitxak} \textsuperscript{(6)} eta \VUgras{urdaia}
\textsuperscript{(7)} saltzen baititu, \VUgras{gurin eta gazta saltzailearen}
\textsuperscript{(8)} ondoan dago, zeinaren mahaia \VUgras{gazta holandarrez}
\textsuperscript{(9)} beteta baitago, azal gorrikoak, \VUgras{camembertez}
\textsuperscript{(10)}, gruyere-\VUgras{gurpil} handiz \textsuperscript{(11)} eta
\VUgras{gurin-puska} freskoz \textsuperscript{(12)}.

%%K t15-03-1 p
3. --- Haren aurrean \VUgras{barazki-saltzaileak} \textsuperscript{(13)} bere
bezeroei eskaintzen dizkie \VUgras{tomate} ederrak \textsuperscript{(14)},
\VUgras{azalorea} \textsuperscript{(15)}, \VUgras{uraza} \textsuperscript{(16)},
\VUgras{aza} \textsuperscript{(17)}, \VUgras{kalabazak}
\textsuperscript{(19)}, \VUgras{meloiak} \textsuperscript{(18)} eta
\VUgras{perretxikoak} \textsuperscript{(20)} ere. Baina \VUgras{garagardogileak},
zeinak bere \VUgras{banaketa-gurdia} \textsuperscript{(21)} gelditu baitu ---
\VUgras{garagardo-barrikaz} \textsuperscript{(22)} kargatua --- haren mahaiaren
aurrean bertan, bezeroei hurbiltzea eragozten die. Baratzezainak
\VUgras{orburu} saski bat \textsuperscript{(23)} ekartzen dio.

%%K t15-tit-2 sub
\VUsaut{4.0mm}
\VUpk{10.2pt}{40}{Saldu eta erosi.}
\VUsaut{3.0mm}

%%K t15-04-1 p
4. --- Merkatuan mugimendu handia dago, \VUgras{salerosketaren}
\textsuperscript{(24)} ordua iritsi baita. \VUgras{Etxekoandreek}
\textsuperscript{(26)}, erosketak egitera etortzen direnek, bidean gelditzen dira
\VUgras{hestebete-saltzailearen} \textsuperscript{(25)} mahaiaren aurrean
\VUgras{tripakiak} edo \VUgras{txahal-burua} erosteko.

%%K t15-05-1 p
5. --- \VUgras{Sukaldari} lodi bat \textsuperscript{(27)} \VUgras{hegaluze}
eder baten prezioa eztabaidatzen ari da \VUgras{arrain-saltzailearekin}
\textsuperscript{(28)}, zeinak nahi bezala igotzen baitu prezioa.
\VUgras{Aingira, platuxa, amuarrain} eta \VUgras{karpa} asko ekarri diote. Haren
\VUgras{bakailaoa} eta haren \VUgras{muskuiluak} oso freskoak dira.

%%K t15-tit-3 sub
\VUsaut{4.0mm}
\VUpk{10.2pt}{40}{Merkea eta garestia.}
\VUsaut{3.0mm}

%%K t15-06-1 p
6. --- \VUgras{Hegazti-saltzailea} \textsuperscript{(29)}, haren ondoan kokatua,
oso zintzoa da. \VUgras{Indioilarra}, \VUgras{antzara}, \VUgras{ahatea} edo
\VUgras{oilo} soila saltzean, ez du prezio bidezkoa baino gehiago eskatzen.

%%K t15-07-1 p
7. --- Plazaren izkinan, lehen solairuan, \VUgras{oihal-saltzaile} batek
\textsuperscript{(30)} ireki du denda berriki, non bezeroek beti aurkituko
baitute \VUgras{mahai-zapi}, \VUgras{serbileta, mantal} eta \VUgras{sukaldeko
eskuoihal} aukera oso ona. Solairu berean \VUgras{labangileak}
\textsuperscript{(31)} jarri du bere tailerra. Merkatarien \VUgras{labanak}
zorrozten ditu eta baratzezainei \VUgras{kimatzeko labanak} egiten dizkie
\VUgras{altzairurik} gogorrenaz.

%%K t15-tit-4 sub
\VUsaut{4.0mm}
\VUpk{10.2pt}{40}{Janari-denda.}
\VUsaut{3.0mm}

%%K t15-08-1 p
8. --- Haren tailerraren azpian janari-denda handi bat ireki da duela gutxi. Ez da
jada iraganeko \VUgras{janari-dendatxo} \textsuperscript{(32)} xume eta apala, non
eguneroko gauzak baino ez baitziren saltzen --- \VUgras{tea}
\textsuperscript{(33)}, \VUgras{txokolatea} \textsuperscript{(34)},
\VUgras{azukre-koskorra} \textsuperscript{(37)} eta \VUgras{xaboia}
\textsuperscript{(38)}. Hemen dena saltzen dute: \VUgras{galleta kurruskariak},
\VUgras{kontserbak} \textsuperscript{(35)}, \VUgras{likoreak}
\textsuperscript{(36)}, \VUgras{rona, koinaka, kirsch-a, anisa} eta
\VUgras{curaçao-a}. Ehiza ere erraz lortzen da hemen --- \VUgras{erbia}
\textsuperscript{(39)}, \VUgras{birigarroak} \textsuperscript{(40)},
\VUgras{eperrak} \textsuperscript{(41)}, \VUgras{galeperrak}
\textsuperscript{(42)}, \VUgras{basahatea} \textsuperscript{(43)} edo
\VUgras{faisaia} \textsuperscript{(44)} --- eta baita tropikoetako fruituak ere,
\VUgras{bananak} \textsuperscript{(46)} eta \VUgras{ananak} bezala.

%%K t15-09-1 p
9. --- \VUgras{Saltzaile} oso adeitsuak \textsuperscript{(45)} eskatzen den
guztia emateko prest dago: \VUgras{marmelada} \textsuperscript{(47)},
andere-mahatsezkoa edo irasagarrezkoa, \VUgras{txerri-gantza}
\textsuperscript{(48)}, \VUgras{ozpin-luzokerrak} \textsuperscript{(49)} edo
\VUgras{sardina} gaziak \textsuperscript{(50)}.

%%K t15-09-2 p
Oso erosoa da \VUgras{erosleentzat} \textsuperscript{(51)}, zeinek gauzarik
arruntenen ondoan aurkitzen baitituzte goiztiarrik bakanenak eta fruiturik
goxoenak: \VUgras{udare} zukutsuak \textsuperscript{(52)}, \VUgras{arana}
\textsuperscript{(53)}, \VUgras{mahatsa} \textsuperscript{(54)},
\VUgras{mertxikak} \textsuperscript{(55)}, \VUgras{arbendola}
\textsuperscript{(56)} eta \VUgras{intxaurrak} \textsuperscript{(57)}.

%%K t15-tit-5 sub
\VUsaut{4.0mm}
\VUpk{10.2pt}{40}{Erosleak.}
\VUsaut{3.0mm}

%%K t15-10-1 p
10. --- \VUgras{Arroza} \textsuperscript{(58)} eta \VUgras{dilista}
\textsuperscript{(59)} \VUgras{sardinzar} ketuen \textsuperscript{(60)} kutxaren
ondoan daude; \VUgras{patata} \textsuperscript{(61)} eta \VUgras{indaba}
\textsuperscript{(63)} \VUgras{sagarren} \textsuperscript{(62)},
\VUgras{pikuen} \textsuperscript{(64)}, \VUgras{aran lehorren}
\textsuperscript{(65)} eta \VUgras{hurren} \textsuperscript{(66)} auzokoak dira.
\VUgras{Arrautza} freskoak \textsuperscript{(67)} eta \VUgras{olibak}
\textsuperscript{(68)} ere badaude denda honetan, eta hala \VUgras{saskiak}
\textsuperscript{(70)} eta \VUgras{sareak} \textsuperscript{(73)} oso azkar
betetzen dira.

%%K t15-11-1 p
11. --- Etxe bikain honetako \VUgras{kafea} \textsuperscript{(72)} merezimenduz da
ospetsua \VUgras{neskameen} \textsuperscript{(69)} eta sukaldarien artean,
\VUgras{janari-saltzaile} zaharrak \textsuperscript{(71)} berak erretzen baitu
ardura handienarekin.

%%K t15-tit-6 sub
\VUsaut{4.0mm}
\VUpk{10.2pt}{40}{Zer ikus daitekeen hemen.}
\VUsaut{3.0mm}

%%K t15-12-1 p
12. --- Ez dira denak janarien bila etortzen merkatura. Teilapera daraman
kaletxoaren mugimendu pintoreskoan jende mota guztia ikusten da.
\VUgras{Hiltegiko langile} onbera baten \textsuperscript{(75)} ondoan --- zeinak
\VUgras{eskorga} bat \textsuperscript{(74)} bultzatzen baitu aurretik --- hona
hemen bi soldadu oporretan, beren uniforme distiratsuez oso harro:
\VUgras{husarra} \textsuperscript{(76)} bere \VUgras{dolman} urdinarekin eta
\VUgras{txapel lumadunarekin}, eta \VUgras{zuaboen kaporala} galtza zabalekin eta
\VUgras{feza} buruan.

%%K t15-13-1 p
13. --- \VUgras{Okinak} \textsuperscript{(80)}, \VUgras{okindegiaren}
\textsuperscript{(78)} atalasean dagoenak, bere \VUgras{ogiaren}
\textsuperscript{(79)} orea oratu berri du eta labetik atera ditu
\VUgras{opilak} \textsuperscript{(81)}, \VUgras{txigortxoak}
\textsuperscript{(82)} eta \VUgras{ogi biribilak} \textsuperscript{(83)},
erakusleihoan daudenak.

%%K t15-14-1 p
14. --- Okindegiaren gainean \VUgras{jostun} trebe bat \textsuperscript{(84)} bizi
da, zeinaren etxean \VUgras{brodatzaile} batzuek \textsuperscript{(85)} lan egiten
baitute. Haren ondoan \VUgras{garbitzaile eta lisatzaile} bat
\textsuperscript{(86)} bizi da, zeinak \VUgras{oihalak} \textsuperscript{(88)}
almidoiztatzen baititu, garbitu eta lehortu ondoren, eta
\VUgras{lisaburdinaz} \textsuperscript{(87)} lisatzen.

%%K t15-tit-7 sub
\VUsaut{4.0mm}
\VUpk{10.2pt}{40}{Haragia, arraina eta barazkiak.}
\VUsaut{3.0mm}

%%K t15-15-1 p
15. --- Beheko solairuko \VUgras{harategian} \textsuperscript{(89)} haragi mota
guztiak saltzen dituzte --- \VUgras{arkume-haragia} \textsuperscript{(90)},
\VUgras{behi-haragia} \textsuperscript{(94)} edo \VUgras{txahal-haragia}
\textsuperscript{(92)} --- \VUgras{arkume-hanka, xerra, errekari} eta
\VUgras{txuleta} gisa. \VUgras{Harakinaren morroiak}
\textsuperscript{(93)} haragi hori guztia zatitzen du eta \VUgras{hezur-muina}
\textsuperscript{(96)} bereizten. Harategiaren atean \VUgras{ostra-saltzailea}
\textsuperscript{(97)} dago, zeinak \VUgras{ostrak} \textsuperscript{(98)}
saltzen baititu. Eskatzen bazaio, irekitzen ditu.

%%K t15-16-1 p
16. --- Merkatari horiek guztiek patentea edo lekuaren zerga ordaintzen dute;
horregatik \VUgras{udaltzainak} \textsuperscript{(100)} errukirik gabe uxatzen
ditu \VUgras{barazki-saltzaile ibiltari} gaixoak \textsuperscript{(99)}, zeinek
lehia egiten baitiete beren orgatxoetan \VUgras{azenarioa, errefautxoa, arbia,
porrua} edo \VUgras{zainzuri sortak, ilar freskoak, ziazerbak, mihiluak,
berroa} eta \VUgras{perrexila} banatuz.

%%K t15-tit-8 sub
\VUsaut{4.0mm}
\VUpk{10.2pt}{40}{Kontuz udaltzainarekin!}
\VUsaut{3.0mm}

%%K t15-17-1 p
17. --- \VUgras{Baratxuria} \textsuperscript{(101)} eta \VUgras{tipula} saltzen
dituen mutilak ondo daki berari ere ez zaiola zilegi bere salgaia merkatuaren
ondoan saltzea. Korrika hasten da, \VUgras{karramarro}, \VUgras{ibai-karramarro}
eta \VUgras{izkira} saski bat irauliko duen arriskuan, zeina \VUgras{langosta}
eta \VUgras{otarrain} \VUgras{saltzailearena} baita \textsuperscript{(102)}.

%%K t15-18-1 p
18. --- Denak lanpetuta daude eta zarata handia ateratzen dute; baina horrek ez du
eragozten argi entzutea \VUgras{beirategile} ibiltariaren \textsuperscript{(103)}
deiadarra eta \VUgras{ikatz-saltzailearen} \textsuperscript{(104)} oihua, zeinak
bere orgatxoa une batez utzi baitu \VUgras{ikatz-zaku} bat
\textsuperscript{(105)} eramateko.
\VUsaut{6.0mm}
\VUfilet{22mm}
